% !TEX program = xelatex
% Lernplatt - by Can Siebert
% Kurs 22 (Bo 143) - Tag 9 - Loesungsheft
% Performance-Marketing: Anzeigen, Conversion & Landing Pages

\documentclass[11pt,a4paper,oneside]{article}

\usepackage{polyglossia}
\setdefaultlanguage{german}

\usepackage[a4paper,margin=2.5cm]{geometry}

\usepackage{fontspec}
\defaultfontfeatures{Ligatures=TeX}

\IfFontExistsTF{Charter}{\setmainfont{Charter}}{%
  \setmainfont{Noto Serif}}
\IfFontExistsTF{Source Sans 3}{\setsansfont{Source Sans 3}}{%
  \IfFontExistsTF{Source Sans Pro}{\setsansfont{Source Sans Pro}}{%
    \setsansfont{Liberation Sans}}}

\newcommand{\caveatfontfallback}{\itshape}
\IfFontExistsTF{Caveat}{\newfontfamily\caveatfont{Caveat}[Scale=1.15]}{\newcommand\caveatfont{\caveatfontfallback}}

\setmonofont{Liberation Mono}

\usepackage{microtype}
\usepackage{eurosym}
\linespread{1.15}

\usepackage{xcolor}
\definecolor{lpInk}{RGB}{20,28,38}
\definecolor{lpAccent}{RGB}{22,90,140}
\definecolor{lpAccentLight}{RGB}{210,228,240}
\definecolor{lpAmber}{RGB}{222,138,30}
\definecolor{lpAmberLight}{RGB}{255,243,217}
\definecolor{lpAmberBorder}{RGB}{196,118,18}
\definecolor{lpGray}{RGB}{96,104,114}
\definecolor{lpGrayLight}{RGB}{238,240,243}
\definecolor{lpGreen}{RGB}{34,120,72}
\definecolor{lpGreenLight}{RGB}{222,238,228}
\definecolor{lpGreenBorder}{RGB}{24,96,58}
\definecolor{lpL1}{RGB}{34,120,72}
\definecolor{lpL2}{RGB}{22,90,140}
\definecolor{lpL3}{RGB}{170,42,52}

\usepackage[most]{tcolorbox}
\tcbuselibrary{breakable,skins}

\newtcolorbox{infobox}[1][Hinweis]{enhanced, breakable, colback=lpAccentLight, colframe=lpAccent, boxrule=0.6pt, arc=2pt, left=10pt,right=10pt,top=8pt,bottom=8pt, fonttitle=\sffamily\bfseries\color{white}, coltitle=white, title={#1}, attach boxed title to top left={xshift=8pt,yshift=-8pt}, boxed title style={size=small, colback=lpAccent, colframe=lpAccent, boxrule=0.4pt, arc=1pt}, top=14pt}

\newtcolorbox{loesungbox}[1][Musterlösung]{enhanced, breakable, colback=lpGreenLight, colframe=lpGreenBorder, boxrule=0.6pt, arc=2pt, left=10pt,right=10pt,top=8pt,bottom=8pt, fonttitle=\sffamily\bfseries\color{white}, coltitle=white, title={#1}, attach boxed title to top left={xshift=8pt,yshift=-8pt}, boxed title style={size=small, colback=lpGreenBorder, colframe=lpGreenBorder, boxrule=0.4pt, arc=1pt}, top=14pt}

\newtcolorbox{merkbox}[1][Managementhinweis]{enhanced, breakable, colback=lpAmberLight, colframe=lpAmberBorder, boxrule=0.6pt, arc=2pt, left=10pt,right=10pt,top=8pt,bottom=8pt, fonttitle=\sffamily\bfseries\color{white}, coltitle=white, title={#1}, attach boxed title to top left={xshift=8pt,yshift=-8pt}, boxed title style={size=small, colback=lpAmberBorder, colframe=lpAmberBorder, boxrule=0.4pt, arc=1pt}, top=14pt}

\usepackage{enumitem}
\setlist{itemsep=2pt, topsep=4pt, parsep=0pt}
\setlist[itemize,1]{label={\textbullet},leftmargin=1.6em}
\setlist[itemize,2]{label={\textendash},leftmargin=1.4em}
\setlist[enumerate,1]{label=\arabic*., leftmargin=1.8em}

\usepackage{tabularx}
\usepackage{array}
\usepackage{booktabs}
\usepackage{longtable}

\usepackage{fancyhdr}
\usepackage{lastpage}
\pagestyle{fancy}
\fancyhf{}
\renewcommand{\headrulewidth}{0.3pt}
\renewcommand{\footrulewidth}{0pt}
\fancyhead[L]{\footnotesize\sffamily\color{lpGray}Lösungsheft \textperiodcentered{} Kurs 22 \textperiodcentered{} Tag 9}
\fancyhead[R]{\footnotesize\sffamily\color{lpGray}Performance-Marketing}
\fancyfoot[L]{\footnotesize\sffamily\color{lpGray}\textcopyright{} Can Siebert}
\fancyfoot[R]{\footnotesize\sffamily\color{lpGray}\thepage{} / \pageref{LastPage}}

\fancypagestyle{plain}{\fancyhf{}\renewcommand{\headrulewidth}{0pt}\fancyfoot[C]{\footnotesize\sffamily\color{lpGray}\thepage{} / \pageref{LastPage}}}

\usepackage[explicit]{titlesec}
\titleformat{\section}[block]{\sffamily\Large\bfseries\color{lpAccent}}{\thesection.}{0.6em}{#1}[{\vspace{2pt}\color{lpAccent}\titlerule[0.6pt]}]
\titleformat{\subsection}[block]{\sffamily\large\bfseries\color{lpInk}}{\thesubsection}{0.6em}{#1}
\titlespacing*{\section}{0pt}{14pt}{8pt}
\titlespacing*{\subsection}{0pt}{10pt}{4pt}

\usepackage[hidelinks,unicode]{hyperref}

\newcommand{\akzent}[1]{{\caveatfont\color{lpAmber}#1}}
\newcommand{\fachbegriff}[1]{\textbf{#1}}

\newcommand{\niveaubadge}[2]{\tcbox[on line, colback=#1, colframe=#1, coltext=white, boxrule=0pt, arc=2pt, left=5pt,right=5pt,top=1pt,bottom=1pt, nobeforeafter]{\sffamily\bfseries\footnotesize #2}}
\newcommand{\nivLi}{\niveaubadge{lpL1}{L1\,\textbullet\,Wissen}}
\newcommand{\nivLii}{\niveaubadge{lpL2}{L2\,\textbullet\,Anwendung}}
\newcommand{\nivLiii}{\niveaubadge{lpL3}{L3\,\textbullet\,Management}}

\newcommand{\zeit}[1]{\tcbox[on line, colback=lpGrayLight, colframe=lpGrayLight, coltext=lpInk, boxrule=0pt, arc=2pt, left=5pt,right=5pt,top=1pt,bottom=1pt, nobeforeafter]{\sffamily\footnotesize\textbf{$\sim$\,#1}}}

\newcommand{\aufgabenkopf}[4]{\par\vspace{6pt}\noindent{\sffamily\Large\bfseries\color{lpAccent}Aufgabe #1 \textendash{} #2}\par\vspace{2pt}\noindent{\color{lpAccent}\rule{\linewidth}{0.6pt}}\par\vspace{4pt}\noindent #3 \hfill \zeit{#4}\par\vspace{6pt}}

\newcommand{\ytquery}[1]{\par\vspace{4pt}\noindent{\footnotesize\sffamily\color{lpGray}\textbf{YouTube-Suchquery:} \texttt{#1}}\par}

\begin{document}

\begin{titlepage}
  \pagestyle{empty}
  \vspace*{1.5cm}
  \noindent{\sffamily\footnotesize\color{lpGray}LERNPLATT \textperiodcentered{} BY CAN SIEBERT}\\[2pt]
  \noindent{\color{lpAccent}\rule{\linewidth}{1.2pt}}\\[4pt]
  \noindent{\sffamily\footnotesize\color{lpGray}LÖSUNGSHEFT \textperiodcentered{} MODUL 5 \textperiodcentered{} TAG 9 \textperiodcentered{} KURS 22 (Bo 143) \textperiodcentered{} 10.\,Juni 2026}

  \vspace{3cm}
  {\sffamily\fontsize{30}{34}\selectfont\bfseries\color{lpInk}Lösungsheft\\Tag 9\par}

  \vspace{0.7cm}
  {\sffamily\large\color{lpGray}Musterlösungen zu den elf Aufgaben\\des Aufgabenhefts \textendash{} Performance-\\Marketing im Vertrieb.\par}

  \vspace{1.5cm}
  {\caveatfont\LARGE\color{lpGreen}Musterlösungen \textperiodcentered{} nicht die einzige Antwort}\par

  \vfill

  \noindent\begin{minipage}{0.55\linewidth}
    {\sffamily\footnotesize\color{lpGray}Hinweis:}\\
    {\sffamily\small\color{lpInk}Musterlösungen sind Diskussionsbasis.}\\[10pt]
    {\sffamily\footnotesize\color{lpGray}Begleitend:}\\
    {\sffamily\small\color{lpInk}Aufgabenheft \& Lehrskript Tag 9}
  \end{minipage}\hfill
  \begin{minipage}{0.4\linewidth}
    \raggedleft
    {\sffamily\footnotesize\color{lpGray}Dozent:}\\
    {\sffamily\small\color{lpInk}Can Siebert}\\[10pt]
    {\sffamily\footnotesize\color{lpGray}Niveau-Mix:}\\
    {\sffamily\small\color{lpInk}3\,L1 \textperiodcentered{} 5\,L2 \textperiodcentered{} 3\,L3}
  \end{minipage}

  \vspace{0.5cm}
  \noindent{\color{lpAccent}\rule{\linewidth}{0.6pt}}
\end{titlepage}

\section*{Hinweise zu den Musterlösungen}
\thispagestyle{plain}

\noindent Eine mögliche Lösung. Mehrere Begründungswege sind legitim, solange sie:

\begin{itemize}
  \item den Begriffsapparat von Tag 9 nutzen (Funnel, KPIs, Anzeigenkanäle, Landing Page),
  \item Annahmen sichtbar machen,
  \item Zielkonflikte (Leadmenge vs.\ Leadqualität, Kosten vs.\ Conversion) benennen,
  \item zu einer klaren Investitionsentscheidung führen.
\end{itemize}

\clearpage

\section*{Niveau L1 \textendash{} Wissen \& Reproduktion}
\addcontentsline{toc}{section}{L1 -- Wissen}

% --- AUFGABE 1 -----------------------------------------------------------
% LERNPLATT_HILFSVIDEOS
\section*{Hilfsvideos zu diesem Tag}
\addcontentsline{toc}{section}{Hilfsvideos zu diesem Tag}
\thispagestyle{plain}

\noindent Die folgenden YouTube-Erkl\"arvideos vermitteln die Inhalte, die Sie zur Bearbeitung der Aufgaben ben\"otigen. Klicken Sie auf den Titel, um das Video direkt zu \"offnen; die vollst\"andige URL steht in Klammern.

\begin{description}[style=nextline,leftmargin=18pt,labelindent=0pt,itemsep=8pt]
  \item[1.~Performance-Marketing als messbarer Vertriebshebel] \href{https://youtu.be/756dea00Wmg}{\textcolor{lpAccent}{https://youtu.be/756dea00Wmg}}\\[2pt]
  {\footnotesize\color{lpGray}(URL: \nolinkurl{https://youtu.be/756dea00Wmg})}
  \item[2.~Landing Page Optimierung: Conversion-Hebel in der Praxis] \href{https://youtu.be/ZOE-CCudzZM}{\textcolor{lpAccent}{https://youtu.be/ZOE-CCudzZM}}\\[2pt]
  {\footnotesize\color{lpGray}(URL: \nolinkurl{https://youtu.be/ZOE-CCudzZM})}
  \item[3.~Sales Funnel: vom Klick zum Abschluss] \href{https://youtu.be/ZBEnOEcXtcs}{\textcolor{lpAccent}{https://youtu.be/ZBEnOEcXtcs}}\\[2pt]
  {\footnotesize\color{lpGray}(URL: \nolinkurl{https://youtu.be/ZBEnOEcXtcs})}
  \item[4.~Kennzahlen lesen: KPI-Architektur entlang des Funnels] \href{https://youtu.be/hmpDH5K5jRU}{\textcolor{lpAccent}{https://youtu.be/hmpDH5K5jRU}}\\[2pt]
  {\footnotesize\color{lpGray}(URL: \nolinkurl{https://youtu.be/hmpDH5K5jRU})}
\end{description}
\vspace{0.6em}
\noindent{\footnotesize\color{lpGray}\textit{Tipp:} \"Offnen Sie das Video, lassen Sie es im Hintergrund laufen und arbeiten Sie parallel an der Aufgabe.}

\clearpage

\aufgabenkopf{1}{Reichweite \textendash{} Traffic \textendash{} Leads \textendash{} Conversion \textendash{} Umsatz}{\nivLi}{6\,min}

\ytquery{Performance Marketing Funnel Reichweite Traffic Lead Conversion Umsatz}

\begin{loesungbox}
\textbf{1. Definitionen:}
\begin{itemize}
  \item \emph{Reichweite:} Wie viele Menschen werden potenziell erreicht (Impressionen).
  \item \emph{Traffic:} Wie viele tatsächlich auf eine Landing Page kommen (Klicks).
  \item \emph{Leads:} Wie viele Kontaktinformationen hinterlassen (Formular).
  \item \emph{Conversion:} Wie viele dieser Leads tatsächlich zum Geschäft werden (SQL, Angebot, Abschluss).
  \item \emph{Umsatz:} Wirtschaftliches Endergebnis nach allen Stufen.
\end{itemize}

\textbf{2. Gemessen vs.\ messen müssen:}
Im Mittelstand wird am häufigsten auf der Ebene \emph{Reichweite / Traffic / Leads} gemessen \textendash{} weil sie operativ sichtbar und im Marketing-Tool reportiert sind. Gemessen werden müsste primär auf \emph{Conversion und Umsatz / Deckungsbeitrag} \textendash{} weil nur diese Ebenen wirtschaftliche Wirkung zeigen.

\textbf{3. Warum Reichweite in B2B nicht funktioniert:}
B2B-Märkte sind klein und spezifisch. Eine ``Reichweiten''-Optimierung erreicht eine breite, undefinierte Masse \textendash{} davon ist nur ein Bruchteil im richtigen Buying Center, mit echtem Bedarf, in der richtigen Phase. Streuverluste sind in B2B oft 95\,\%+. Wer auf Reichweite optimiert, kauft Aktivität ohne Vertriebsrelevanz.

\begin{merkbox}[Managementhinweis]
\emph{Reichweite ist eine B2C-KPI. Im B2B führt sie zu Leadmüll.} Die einzige akzeptable Reichweiten-Verwendung in B2B ist als \emph{Frühindikator für Markenwirkung} \textendash{} nicht als Erfolgsmaß einer Kampagne.
\end{merkbox}
\end{loesungbox}

\clearpage

% --- AUFGABE 2 -----------------------------------------------------------
\aufgabenkopf{2}{Zehn Performance-Marketing-KPIs}{\nivLi}{8\,min}

\ytquery{Performance Marketing KPI CPC CPL CAC ROAS Conversion Funnel}

\begin{loesungbox}
\textbf{1. KPI-Zuordnung zu Funnel-Ebenen:}

\smallskip
\begin{tabularx}{\linewidth}{@{}lX@{}}
\toprule
\textbf{Ebene} & \textbf{KPI} \\
\midrule
Anzeige                & Impressionen, Klickrate, CPC \\
\midrule
Klick / Landing        & Conversion Rate (Landing Page) \\
\midrule
Lead                   & CPL, Leadqualität \\
\midrule
Vertrieb               & Abschlussquote, CAC \\
\midrule
Wirtschaftlichkeit     & ROAS, Deckungsbeitrag pro Kampagne \\
\bottomrule
\end{tabularx}

\textbf{2. Wichtigster Frühindikator:}
\emph{Leadqualität} (typischerweise gemessen als \%-Anteil MQL\,$\rightarrow$\,SQL). Sie ist der einzige Frühindikator, der noch \emph{vor} Abschluss zeigt, ob eine Kampagne wirtschaftlich werden kann. Hohe Lead-Mengen mit niedriger Qualität sind eine teure Sackgasse, die Wochen später als schlechter ROAS sichtbar wird.

\textbf{3. Am häufigsten nicht berichtet:}
\emph{Deckungsbeitrag pro Kampagnenquelle}. Wird selten berichtet, weil er Marketing- und Controlling-Daten verbinden muss \textendash{} und das ist organisatorisch unbequem. Folge: Kampagnen, die viele Leads bringen, aber wirtschaftlich Verluste produzieren, laufen monatelang weiter, weil sie ``erfolgreich'' aussehen.

\begin{merkbox}[Managementhinweis]
\emph{Deckungsbeitrag pro Kampagnenquelle ist die wichtigste KPI, die kein Mittelständler berichtet.} Wer diese Zahl monatlich auf den Tisch legt, verändert das Anreizsystem von Marketing strukturell.
\end{merkbox}
\end{loesungbox}

\clearpage

% --- AUFGABE 3 -----------------------------------------------------------
\aufgabenkopf{3}{Anzeigenkanäle nach Suchintention}{\nivLi}{6\,min}

\ytquery{Performance Marketing Google Ads LinkedIn Suchintention Kaufabsicht B2B}

\begin{loesungbox}
\textbf{1. Zuordnung:}

\smallskip
\begin{tabularx}{\linewidth}{@{}lX@{}}
\toprule
\textbf{Stufe} & \textbf{Kanal} \\
\midrule
Problemrecherche    & Display, LinkedIn (Content), Meta (Interessen), Plattformanzeigen, Partnerkampagnen \\
\midrule
Lösungsvergleich    & Google Ads (informative Keywords), LinkedIn Ads (Whitepaper, Webinar), Retargeting \\
\midrule
Kaufabsicht         & Google Ads (kaufnahe Keywords), Retargeting, Plattformanzeigen \\
\bottomrule
\end{tabularx}

\textbf{2. Konkrete Kaufabsicht \textendash{} wirksamster Kanal:}
\emph{Google Ads für kaufnahe Keywords} (``Cybersecurity Beratung Mittelstand'', ``ERP-System B2B kaufen''). Diese Suchen entstehen, wenn ein konkreter Bedarf da ist \textendash{} der Klick führt zu hoher Wahrscheinlichkeit der Conversion, weil die Suchintention bereits Kaufabsicht enthält.

\textbf{3. Für B2B-Buying-Center-Ansprache:}
\emph{LinkedIn Ads} (Targeting nach Funktion, Branche, Unternehmensgröße) und \emph{Partnerkampagnen} (Co-Marketing mit Branchenpartnern, die das Buying Center bereits kennen). Beide ermöglichen präzise Ansprache der relevanten Rollen (Initiator, Influencer, Decider), die im Buying Center entscheiden.
\end{loesungbox}

\clearpage

\section*{Niveau L2 \textendash{} Anwendung \& Transfer}
\addcontentsline{toc}{section}{L2 -- Anwendung}

% --- AUFGABE 4 -----------------------------------------------------------
\aufgabenkopf{4}{SolarOffice GmbH \textendash{} Performance-Marketing-Funnel}{\nivLii}{25\,min}

\ytquery{Performance Marketing B2B PV Solar Gewerbe Zielgruppe Funnel Kampagne}

\begin{loesungbox}
\textbf{1. Vertriebsziel:}
SolarOffice will in 6\,Monaten 80 qualifizierte Beratungstermine generieren, aus denen mindestens 20 Aufträge entstehen sollen. Ziel ist \emph{nicht} ``mehr Leads'', sondern \emph{Wachstum der Pipeline mit Konversionsfähigkeit} \textendash{} bei begrenzter Vertriebskapazität.

\textbf{2. Drei Zielgruppen:}
\begin{itemize}
  \item Eigentümer / GF mittelständischer Produktionsbetriebe mit hohem Stromverbrauch.
  \item Immobilieneigentümer / Hausverwaltungen von Gewerbeimmobilien.
  \item Logistik- und Lagerbetriebe mit großen Dachflächen.
\end{itemize}

\textbf{3. Drei Kundenprobleme für Anzeigen:}
\begin{itemize}
  \item ``Energiekosten verdoppeln sich \textendash{} Sie suchen Planungssicherheit für die nächsten 10\,Jahre.''
  \item ``Sie wollen Fördermittel nutzen, aber haben keine Zeit, sich durch Formulare zu kämpfen.''
  \item ``Sie wollen Nachhaltigkeit zeigen \textendash{} mit klar berechneter Wirtschaftlichkeit, nicht mit Greenwashing.''
\end{itemize}

\textbf{4. Funnel (5 Stufen):}
\begin{enumerate}
  \item LinkedIn-/Google-Anzeige mit Problem-Ansprache.
  \item Landing Page mit Energiekosten-Rechner und Beispielprojekten.
  \item Whitepaper-Download oder Beratungstermin-Anfrage (Formular).
  \item Innendienst kontaktiert qualifizierten Lead innerhalb 24\,h.
  \item Beratungstermin Vor-Ort oder Video.
\end{enumerate}

\textbf{5. Drei Kennzahlen (nicht Klicks):}
\begin{itemize}
  \item Beratungstermine geliefert (Hauptziel-KPI).
  \item Konversion Anfrage $\rightarrow$ Beratungstermin (Lead-Qualität).
  \item Kosten pro Beratungstermin (CPL auf Vertriebs-Ebene).
\end{itemize}
\end{loesungbox}

\clearpage

% --- AUFGABE 5 -----------------------------------------------------------
\aufgabenkopf{5}{SolarOffice \textendash{} Kanäle \& Landing Page}{\nivLii}{30\,min}

\ytquery{Performance Marketing Landing Page B2B Solar Gewerbe Google LinkedIn}

\begin{loesungbox}
\textbf{1. Kanalbewertung:}

\smallskip
\begin{tabularx}{\linewidth}{@{}lccccc@{}}
\toprule
& \textbf{Ziel-G.} & \textbf{Risiko} & \textbf{Lead-Q.} & \textbf{Geschw.} \\
\midrule
A: Google Ads (Suchanfragen)    & 5 & 3 & 4 & 5 \\
B: LinkedIn Ads                  & 5 & 3 & 5 & 4 \\
C: Display                       & 2 & 5 & 1 & 5 \\
D: Retargeting                   & 4 & 2 & 4 & 4 \\
E: Plattform-/Energieportale    & 4 & 3 & 4 & 3 \\
\bottomrule
\end{tabularx}

\smallskip
\emph{Risiko: 5\,=\,hoch, 1\,=\,niedrig.}

\textbf{2. Drei Kanäle:} \textbf{A} (Google Ads), \textbf{B} (LinkedIn Ads), \textbf{D} (Retargeting). \emph{Bewusst nicht:} Display (Streuverluste) und Plattform-Portale in Phase 1 (erst nach Lerneffekt).

\textbf{3. Fünf Schwächen der Landing Page:}
\begin{itemize}
  \item Überschrift ist generisch (``Willkommen bei SolarOffice'') \textendash{} kein Problembezug.
  \item Keine Zielgruppenansprache für Gewerbekunden.
  \item Formular mit 12 Pflichtfeldern \textendash{} massive Abbruchquelle.
  \item Keine Beispielrechnung / kein konkreter Mehrwert.
  \item Keine Referenzen oder Vertrauenselemente.
  \item CTA ``Absenden'' ist nichts-sagend, lädt nicht zur Aktion ein.
\end{itemize}

\textbf{4. Fünf konkrete Verbesserungen:}
\begin{itemize}
  \item Überschrift: ``Photovoltaik für Ihr Gewerbe \textendash{} 30\,\% niedrigere Energiekosten in 12\,Monaten.''
  \item Zielgruppenansprache: ``Für Geschäftsführer und Eigentümer von Produktions-, Lager- und Gewerbeimmobilien.''
  \item Formular auf 4 Felder reduzieren (Name, Firma, Stromverbrauch / Dachfläche, Telefon).
  \item Beispielrechnung: ``2.000\,m\textsuperscript{2} Dach, 1\,Mio.\ kWh/Jahr Verbrauch $\rightarrow$ ca.\ 180.000\,\euro{} Einsparung in 10\,Jahren.''
  \item 3 Referenz-Kundenlogos + 1 Case-Quote: ``Wir haben in 2\,Jahren 320.000\,\euro{} eingespart''.
  \item CTA: ``Kostenlose Wirtschaftlichkeitsanalyse für Ihr Gewerbeobjekt anfragen''.
\end{itemize}

\textbf{5. Budget zuerst Anzeigen oder Landing Page?}
\textbf{Zuerst Landing Page} (mind.\ 60\,\% der ersten 10.000\,\euro). Begründung: Eine gute Anzeige mit schlechter Landing Page verbrennt Klick-Geld. Eine schlechte Anzeige mit guter Landing Page ist immerhin reparabel. Conversion-Hub muss stehen, bevor Klick-Volumen aufgebaut wird.

\begin{merkbox}[Managementhinweis]
\emph{60\,\% der typischen B2B-Mittelstands-Performance-Budgets sollten in die Landing Page fließen \textendash{} 40\,\% in die Anzeigen.} In der Praxis ist das Verhältnis meist umgekehrt \textendash{} und das ist die häufigste Ursache schwacher Performance.
\end{merkbox}
\end{loesungbox}

\clearpage

% --- AUFGABE 6 -----------------------------------------------------------
\aufgabenkopf{6}{Elemente einer wirksamen Landing Page}{\nivLii}{15\,min}

\ytquery{Landing Page B2B Überschrift Nutzenversprechen Vertrauen Conversion}

\begin{loesungbox}
\textbf{1. Acht Elemente und ihre Rolle:}
\begin{itemize}
  \item \emph{Klare Überschrift:} liefert in 3\,Sek.\ Antwort auf ``Bin ich hier richtig?''.
  \item \emph{Nutzenversprechen:} klärt sofort, was der Besucher konkret gewinnt.
  \item \emph{Zielgruppenansprache:} signalisiert ``wir verstehen Ihre Situation''.
  \item \emph{Vertrauenselemente / Referenzen:} reduzieren wahrgenommenes Risiko (Sterne, Logos, Zitate).
  \item \emph{Formular:} so kurz wie möglich, nur Pflichtfelder.
  \item \emph{CTA:} klar, handlungsorientiert, nicht ``Absenden''.
  \item \emph{Visuelle Führung:} Blick lenkt automatisch zum CTA / wichtigen Inhalten.
  \item \emph{Mobile Nutzbarkeit:} 60\,\%+ B2B-Recherche im Mobil \textendash{} ohne mobile Optimierung verloren.
\end{itemize}

\textbf{2. Schwach im Mittelstand:}
\emph{Nutzenversprechen} (oft generisch, keine konkrete Zahl/Aussage) und \emph{Vertrauenselemente} (Referenzen fehlen oder sind zu allgemein gehalten).

\textbf{3. Überschrift + Nutzenversprechen für SolarOffice:}
``\emph{Photovoltaik für Ihr Gewerbe \textendash{} 30\,\% niedrigere Energiekosten in den nächsten 10\,Jahren.} Wir berechnen Wirtschaftlichkeit, Fördermittel und Umsetzungsplan transparent \textendash{} damit Sie wissen, ob sich PV für Ihr Objekt lohnt \textendash{} bevor Sie investieren.''

\textbf{4. Drei Vertrauenselemente im B2B-Gewerbe:}
\begin{itemize}
  \item Konkrete Kunden-Logos aus dem Zielsegment (vergleichbare Größe, Branche).
  \item Zertifikate, Auszeichnungen, Branchenmitgliedschaften.
  \item Konkrete Fallstudie mit Zahlen (\emph{nicht}: ``zufriedene Kunden seit 20\,Jahren'').
\end{itemize}

\begin{merkbox}[Lessons Learned]
\emph{Vertrauen wird im B2B durch \emph{Zahlen}, nicht durch \emph{Adjektive} gestärkt.} ``Hochwertig'' überzeugt niemanden; ``320.000\,\euro{} Einsparung in 2\,Jahren'' schon.
\end{merkbox}
\end{loesungbox}

\clearpage

% --- AUFGABE 7 -----------------------------------------------------------
\aufgabenkopf{7}{DataSecure Solutions \textendash{} Kampagnenstrategie}{\nivLii}{30\,min}

\ytquery{Cybersecurity B2B Performance Marketing LinkedIn Webinar Landing Page Demo}

\begin{loesungbox}
\textbf{1. Zielgruppenanalyse:}
\begin{itemize}
  \item \emph{IT-Leiter:} suchen fachliche Sicherheit, technische Passung, Integrationsfähigkeit.
  \item \emph{GF:} achten auf Risiko, Kosten, Haftung, Geschäftskontinuität.
  \item \emph{Datenschutz / Compliance:} brauchen Nachweise, Prozesse, Auditfähigkeit.
  \item Buying Center mehrstufig $\rightarrow$ Kampagnen müssen unterschiedliche Informationsbedürfnisse bedienen.
\end{itemize}

\textbf{2. Anzeigenplanung 6 Monate:}
\begin{itemize}
  \item Google Ads für akute Nachfrage (Monat 1\,--\,6).
  \item LinkedIn Ads für gezielte Buying-Center-Ansprache (Monat 1\,--\,6, Hauptkanal).
  \item Webinar als Vertrauens- und Qualifikationskanal (alle 6\,--\,8 Wochen).
  \item Retargeting für Webinar-Teilnehmer + Website-Besucher (Monat 2\,--\,6).
  \item \emph{Bewusst nicht:} breite Display-Kampagne.
\end{itemize}

\textbf{3. Priorisierung:}
\begin{itemize}
  \item Priorität 1: \textbf{LinkedIn Ads + Whitepaper} (höchste Leadqualität).
  \item Priorität 2: \textbf{Webinar-Kampagne} (Vertrauensaufbau, fachliche Tiefe).
  \item Priorität 3: \textbf{Google Ads} (kaufnahe Suchen).
  \item Priorität 4: \textbf{Retargeting} (Vertiefung von bestehenden Kontakten).
  \item \emph{Nein:} Display \textendash{} Streuverluste, schwache Leadqualität, hoher Budget-Verbrauch ohne Conversion.
\end{itemize}

\textbf{4. Landing-Page-Struktur:}
\begin{itemize}
  \item Überschrift mit Problembezug: ``Cyberrisiken im Mittelstand erkennen und gezielt reduzieren''.
  \item Kurzes Nutzenversprechen für Entscheider.
  \item Zielgruppenspezifische Abschnitte (IT / GF / Compliance).
  \item Vertrauenselemente: Referenzen, Zertifizierungen (ISO-27001 etc.), Kundenstimmen.
  \item Kurzer Beratungs-/Demo-Prozess (3 Schritte).
  \item CTA: ``Kostenlosen Sicherheitscheck anfragen'' oder ``Demo-Termin vereinbaren''.
  \item Formular mit wenigen Pflichtfeldern + optionale Bedarfsabfrage.
  \item FAQ zu Aufwand, Datenschutz, Ablauf, Kostenrahmen.
\end{itemize}

\textbf{5. Conversion-Ziele + KPIs:}
\begin{itemize}
  \item Klickrate je Kanal/Zielgruppe.
  \item Conversion Rate Landing Page (Ziel: $\geq$\,4\,\%).
  \item Kosten pro Lead (CPL).
  \item Anteil qualifizierter Leads ($\geq$\,40\,\%).
  \item Kosten pro qualifiziertem Lead (CPL-Q).
  \item Demo-Rate ($\geq$\,30\,\% der SQL).
  \item Angebots- und Abschlussquote.
  \item CAC; Deckungsbeitrag je Kampagnenquelle.
\end{itemize}
\end{loesungbox}

\clearpage

% --- AUFGABE 8 -----------------------------------------------------------
\aufgabenkopf{8}{Funnel-Verluste analysieren}{\nivLii}{20\,min}

\ytquery{Performance Marketing Funnel Conversion Verluste Optimierung B2B}

\begin{loesungbox}
\textbf{1. Drei größte Conversion-Verluste:}

\smallskip
\begin{tabularx}{\linewidth}{@{}lXc@{}}
\toprule
\textbf{Übergang} & \textbf{Conversion} & \textbf{Absolut} \\
\midrule
Klicks $\rightarrow$ Landing Page-Besuche (22.000 von 24.000) & 91,7\,\% & \textendash{}2.000 (Tracking) \\
\midrule
Landing Page $\rightarrow$ Formular-Aufruf (4.400 von 22.000) & 20\,\% & \textendash{}17.600 \\
\midrule
Formular-Aufruf $\rightarrow$ Absendung (880 von 4.400)     & 20\,\% & \textendash{}3.520 \\
\midrule
MQL $\rightarrow$ SAL (132 von 660)                        & 20\,\% & \textendash{}528 \\
\midrule
SAL $\rightarrow$ SQL (79 von 132)                         & 60\,\% & \textendash{}53 \\
\bottomrule
\end{tabularx}

\smallskip
\textbf{Hauptverluste:} (a)~Landing Page $\rightarrow$ Formular-Aufruf (\textendash{}17.600 absolut, 80\,\% Verlust), (b)~Formular-Aufruf $\rightarrow$ Absendung (\textendash{}3.520, 80\,\% Verlust), (c)~MQL $\rightarrow$ SAL (\textendash{}528, 80\,\% Verlust).

\textbf{2. Verantwortlichkeiten:}
\begin{itemize}
  \item \emph{Landing Page $\rightarrow$ Formular:} primär \textbf{Landing-Page-Problem} (Nutzenversprechen unklar, Vertrauen schwach, Pfad nicht klar).
  \item \emph{Formular-Aufruf $\rightarrow$ Absendung:} \textbf{Formular-Problem} (zu viele Felder, zu komplex, kein klarer CTA).
  \item \emph{MQL $\rightarrow$ SAL:} \textbf{Lead-Qualifizierungs-/Anzeigen-Problem} (Anzeigen ziehen falsche Personen an; Lead-Definition zu lax).
\end{itemize}

\textbf{3. Kosten pro SQL bei 75.000\,\euro{} Budget:}
75.000\,\euro{} : 79 SQL = \textbf{949\,\euro{} pro SQL}. Wenn ein Abschluss durchschnittlich 50.000\,\euro{} Investitionssumme bringt (B2B-typisch für Cybersecurity), ist das wirtschaftlich vertretbar \textendash{} bei 6 Abschlüssen aus 79 SQL ergibt das einen CAC von 12.500\,\euro.

\textbf{4. Drei Optimierungsmaßnahmen:}
\begin{itemize}
  \item \emph{Landing Page $\rightarrow$ Formular:} Sichtbares Nutzenversprechen, klare CTA-Buttons mehrfach, stärkere Vertrauenselemente (Logos, Zitate, Zahlen). Hypothese: Conversion auf 30\,\%+.
  \item \emph{Formular-Aufruf $\rightarrow$ Absendung:} Felder von 12 auf 4 reduzieren, klarer CTA (``Demo-Termin in 5\,Min.\ vereinbaren''), Trust-Hinweis (``DSGVO-konform''). Hypothese: Conversion auf 35\,\%+.
  \item \emph{MQL $\rightarrow$ SAL:} Lead-Definition schärfen (Pflichtfeld ``Unternehmensgröße'', ``Rolle''), Anzeigen-Targeting verfeinern, Lead Scoring einführen. Hypothese: Conversion auf 35\,\%+.
\end{itemize}

\begin{merkbox}[Managementhinweis]
\emph{Funnel-Engpässe gewinnen Sie immer am untersten Verlust-Knoten.} Wer mehr Anzeigen-Budget einsetzt, ohne den größten Verlust zu fixen, multipliziert nur das Problem.
\end{merkbox}
\end{loesungbox}

\clearpage

\section*{Niveau L3 \textendash{} Management \& Entscheidungsarchitektur}
\addcontentsline{toc}{section}{L3 -- Management}

% --- AUFGABE 9 -----------------------------------------------------------
\aufgabenkopf{9}{DataSecure Solutions \textendash{} 6-Monats-Kampagnenentscheidung}{\nivLiii}{45\,min}

\ytquery{DataSecure Performance Marketing CRO Budget Kanal Lead Qualifizierung Senior}

\begin{loesungbox}
\textbf{1. Priorisierte Kanalentscheidung (4 Kanäle):}
\textbf{LinkedIn Ads + Whitepaper}, \textbf{Webinar-Kampagne}, \textbf{Google Ads (kaufnahe Keywords)}, \textbf{Retargeting}. \emph{Bewusst nicht: breite Display-Kampagne.}

\textbf{2. Budgetverteilung (150.000\,\euro):}
\begin{itemize}
  \item 45.000\,\euro{} \textendash{} LinkedIn Ads (Hauptkanal).
  \item 30.000\,\euro{} \textendash{} Google Ads (kaufnahe).
  \item 25.000\,\euro{} \textendash{} Webinar-Konzept, Bewerbung, Durchführung (3\,--\,4 Webinare).
  \item 20.000\,\euro{} \textendash{} Landing Page, Tracking, Conversion-Optimierung.
  \item 15.000\,\euro{} \textendash{} Whitepaper, Fachcontent, Referenzen.
  \item 10.000\,\euro{} \textendash{} Retargeting.
  \item 5.000\,\euro{} \textendash{} A/B-Tests + Reporting-Setup.
\end{itemize}

\textbf{3. Lead-Qualifizierungs-Logik:}
\begin{itemize}
  \item \emph{MQL:} Profilpassung (Branche, Unternehmensgröße, Rolle); Beispiel: IT-Leiter oder GF eines Unternehmens 100\,--\,1.000 MA.
  \item \emph{SAL:} Innendienst hat innerhalb von 2\,Werktagen Erstkontakt + Bedarf bestätigt.
  \item \emph{SQL:} Bedarf + Budget + Zeitfenster + Entscheider identifiziert.
  \item \emph{SLA:} MQL $\rightarrow$ Innendienst innerhalb 24\,h; SAL $\rightarrow$ Außendienst innerhalb 5\,Werktagen.
\end{itemize}

\textbf{4. Sechs KPIs entlang Funnel:}
\begin{itemize}
  \item Klickrate je Kanal.
  \item Conversion Rate Landing Page.
  \item CPL.
  \item Anteil qualifizierter Leads (\%).
  \item Kosten pro SQL.
  \item ROAS / Deckungsbeitrag pro Kampagnenquelle.
\end{itemize}

\textbf{5. Roadmap (3 Phasen je 2 Monate):}
\begin{itemize}
  \item \emph{Monat 1\,--\,2 \textendash{} Foundation:} Landing Page optimiert, Tracking-Setup, erste LinkedIn-Tests, Google-Setup, 1 Webinar.
  \item \emph{Monat 3\,--\,4 \textendash{} Skalierung:} erfolgreiche Kanäle skalieren, 2 weitere Webinare, Lead-Scoring live, A/B-Optimierung Landing Page.
  \item \emph{Monat 5\,--\,6 \textendash{} Optimierung:} Retargeting voll im Einsatz, Bottlenecks beseitigt, Skalierungsentscheidung für Folgejahr vorbereiten.
\end{itemize}

\textbf{6. Entscheidung gegen reichweitenstarken Kanal:}
DataSecure verzichtet bewusst auf \emph{breite Display-Kampagnen}, obwohl sie kurzfristig die meiste Reichweite und die niedrigsten CPCs liefern würden. Senior-Level-Begründung: (a)~Vertriebsteam ist mit 8 Personen begrenzt \textendash{} jeder unqualifizierte Kontakt blockiert Zeit für qualifizierte Gespräche. (b)~Cybersecurity ist Vertrauenssache \textendash{} eine billige Display-Anzeige beschädigt die Premium-Beratungs-Positionierung. (c)~Die zu erwartende Conversion ist so niedrig, dass der nominale CPC zwar günstig wirkt, der tatsächliche Cost per SQL aber das Doppelte / Dreifache der LinkedIn-Kosten beträgt.

\textbf{7. Entscheidung unter Unsicherheit:}
Investition von 25.000\,\euro{} in Webinare \emph{ohne} klaren ROI-Beweis. Annahme: \emph{Webinare sind für vertrauenssensible B2B-Themen ein überlegener Lead-Qualifikationskanal, weil sie fachliche Tiefe demonstrieren und Buying-Center-Teilnehmer parallel mitnehmen.} Lernindikator: Webinar-Lead-Qualität (Anteil SQL) $\geq$\,40\,\% bis Monat 4. Wenn nicht: Re-Allokation zu LinkedIn + Google.

\begin{merkbox}[Senior-Level-Hinweis]
\emph{Performance-Marketing-Erfolg im B2B wird nicht in CPC gemessen, sondern in Cost per SQL.} CPC-Optimierung führt im B2B fast immer zu Lead-Müll. Senior-Level-Disziplin: \emph{lieber teurere Klicks, die zu qualifizierten Leads werden, als billige Klicks, die den Vertrieb mit Lead-Müll überlasten.}
\end{merkbox}
\end{loesungbox}

\clearpage

% --- AUFGABE 10 ----------------------------------------------------------
\aufgabenkopf{10}{Leadmenge vs.\ Leadqualität}{\nivLiii}{30\,min}

\ytquery{Lead Menge Qualität Performance Marketing CAC ROAS Konflikt B2B}

\begin{loesungbox}
\textbf{1. Konflikt aus Marketing-Sicht:}
Der Performance-Marketing-Manager wird typischerweise an \emph{Lead-Menge und CPL} gemessen. Optimiert er auf diese KPIs, schaltet er Anzeigen mit weiter Zielgruppendefinition, einfachen Lead-Magneten und niedrigschwelligen Formularen. Resultat: viele Leads, niedrige CPL \textendash{} und ein frustrierter Vertrieb, weil 80\,\% der Leads ohne Bedarf sind. Optimiert er auf Qualität, sinkt die Lead-Menge, der CPL steigt \textendash{} und sein Reporting sieht ``schlechter'' aus. Strukturproblem: Anreiz steht gegen Wirkung.

\textbf{2. Gemeinsame Erfolgs-KPI:}
\emph{Cost per SQL} und \emph{Deckungsbeitrag pro Kampagnenquelle}. Beide verbinden Marketing-Output (Lead-Volumen + Lead-Qualifikation) und Vertriebs-Erfolg (Abschlussquote). Marketing kann nicht mehr durch ``mehr Leads'' glänzen, weil schlechte Leads nicht zu SQL werden. Vertrieb kann nicht ``schlechte Leads schuld'' sagen, weil ihre Conversion in der KPI sichtbar ist.

\textbf{3. Anreizgestaltung:}
\begin{itemize}
  \item Marketing-Boni teilweise an Anzahl SQL (nicht MQL) koppeln.
  \item Quartalsweise Marketing-Review mit Vertriebsleitung zur SQL-Qualität.
  \item Gemeinsames Dashboard, das Marketing- und Vertriebs-Daten verbindet.
\end{itemize}

\textbf{4. Lead-Müll-Entlarvung:}
\emph{Anteil MQL, die zu SQL werden.} Bei Kampagnen mit Lead-Müll fällt dieser Anteil auf $<$\,15\,\%. Bei sauberer Kampagne liegt er bei 25\,--\,40\,\%. Diese Zahl ist innerhalb von Wochen sichtbar und entlarvt sofort.

\textbf{5. Kanal-Anfälligkeit:}
\begin{itemize}
  \item \emph{Anfällig für Müll:} Display, breite Meta Ads, generische Plattform-Anzeigen mit nicht-zielgerichteten Lead-Magneten.
  \item \emph{Gut für Qualität:} LinkedIn (Targeting nach Rolle), Webinare (selbst-qualifizierender Effekt), Google Ads (kaufnahe), Partnerkampagnen (vertrauensvorqualifiziert).
\end{itemize}

\begin{merkbox}[Lessons Learned]
\emph{Anreize formen Verhalten. Wer Marketing auf Lead-Menge incentiviert, bekommt Lead-Menge \textendash{} mit allem, was das im B2B bedeutet.} Strukturelle Lösung: KPI-System ändern, nicht ``mehr Workshops''.
\end{merkbox}
\end{loesungbox}

\clearpage

% --- AUFGABE 11 ----------------------------------------------------------
\aufgabenkopf{11}{Transferprojekt \textendash{} Performance-Marketing-Architektur}{\nivLiii}{45\,min}

\ytquery{Performance Marketing Architektur B2B CRO Conversion Lead Qualität Governance}

\begin{loesungbox}
\emph{Musterlösung als Entscheidungsarchitektur \textendash{} andere Konfigurationen möglich.}

\smallskip
\textbf{1. Zielbild:}
Das Unternehmen betreibt in 12\,Monaten eine steuerbare Performance-Marketing-Architektur, die monatlich planbare qualifizierte Leads liefert. Cost per SQL ist die zentrale Steuerungs-KPI. Marketing und Vertrieb arbeiten mit gemeinsamen KPIs. Mindestens 35\,\% des Pipeline-Wertes stammt aus digitalen Quellen. ROAS $\geq$\,3.

\textbf{2. Zielgruppen- und Buying-Center-Analyse:}
\begin{itemize}
  \item \emph{Initiator + User} (Fachverantwortliche): inhaltliche Tiefe wichtig.
  \item \emph{Influencer} (interne/externe Berater): Whitepaper, Webinare.
  \item \emph{Decider} (GF/CFO): ROI, Risiko, Skalierungsfähigkeit.
  \item \emph{Buyer} (Einkauf): Konditionen, Vergleichbarkeit.
\end{itemize}

\textbf{3. Kampagnenstrategie:}
\begin{itemize}
  \item Hauptkanal: LinkedIn Ads (Targeting nach Rolle).
  \item Sekundär: Webinare (Vertrauensaufbau).
  \item Tertiär: Google Ads (kaufnahe Suchen).
  \item Retargeting: Vertiefung bestehender Kontakte.
  \item \emph{Bewusst nicht:} Display, breite Meta Ads.
\end{itemize}

\textbf{4. Botschaftenlogik je Zielgruppe und Funnel-Phase:}
\begin{itemize}
  \item \emph{Problem (Top of Funnel):} Initiator-Schmerz, ``Sind Sie überhaupt sicher?''.
  \item \emph{Lösung (Middle):} Whitepaper, Fallstudien, Webinare.
  \item \emph{Entscheidung (Bottom):} ROI-Rechner, Demo, Referenzen.
\end{itemize}

\textbf{5. Landing-Page-Architektur:}
\begin{itemize}
  \item Pro Kampagne dedizierte Landing Page.
  \item Klare Überschrift mit Problembezug.
  \item Zielgruppen-Differenzierung (sichtbar).
  \item Vertrauenselemente (Logos, Zitate, Zahlen).
  \item Kurzes Formular (4\,--\,5 Felder).
  \item Klare CTA-Hierarchie.
  \item Mobile-first-Design.
  \item FAQ + Tracking + A/B-Tests.
\end{itemize}

\textbf{6. Leadqualifizierungs-Logik:}
\begin{itemize}
  \item MQL: Profilpassung.
  \item SAL: Innendienst-Kontakt erfolgreich, Bedarf bestätigt (binnen 24\,h).
  \item SQL: BANT-Kriterien erfüllt (Budget, Authority, Need, Timeline).
  \item Übergaberegeln im CRM dokumentiert.
\end{itemize}

\textbf{7. KPI-Architektur:}
\begin{itemize}
  \item Anzeige: Klickrate, CPC.
  \item Landing: Conversion Rate.
  \item Lead: CPL, Anteil MQL\,$\rightarrow$\,SQL.
  \item Vertrieb: Abschlussquote, CAC.
  \item Wirtschaftlichkeit: ROAS, Deckungsbeitrag pro Kampagne.
\end{itemize}

\textbf{8. Budget-Logik (Testen / Skalieren / Stoppen / Optimieren):}
\begin{itemize}
  \item \emph{Test-Budget:} max.\ 10\,\% pro neuem Kanal/Format, 4\,Wochen.
  \item \emph{Skalieren:} bei Cost per SQL unter Zielwert + qualitativen Vertriebs-Feedback.
  \item \emph{Stoppen:} bei Cost per SQL 2x Zielwert nach 4\,Wochen \textendash{} kein ``mal sehen ob es besser wird''.
  \item \emph{Optimieren:} A/B-Tests laufen kontinuierlich, mindestens 2 pro Monat.
\end{itemize}

\textbf{9. Governance:}
\begin{itemize}
  \item CRO: Strategie, Budget, Konfliktlösung.
  \item Head of Performance Marketing: operative Steuerung.
  \item Vertriebsleitung: Lead-Übernahme, SLA, Feedback.
  \item Controlling: Audit ROAS, Deckungsbeitrag.
  \item Monatlicher Performance-Marketing-Review-Kreis (CRO, Marketing, Vertrieb, Controlling).
\end{itemize}

\textbf{10. Drei Managemententscheidungen unter Unsicherheit:}
\begin{enumerate}
  \item \emph{Marketing-Anreize teilweise an Cost per SQL koppeln} \textendash{} ohne Beweis, dass das funktioniert. Annahme: Anreiz formt Verhalten besser als Workshops.
  \item \emph{Investition in Webinare} trotz höherer Stückkosten. Annahme: Webinar-Leads konvertieren langfristig signifikant besser.
  \item \emph{Verzicht auf Display} trotz kurzfristigem Reichweiten-Potenzial. Annahme: Lead-Qualität ist im B2B wertvoller als Reichweite.
\end{enumerate}

\textbf{Zusatz \textendash{} Verzicht auf reichweitenstarken Kanal:}
Bewusster Verzicht auf eine breite Meta-Display-Kampagne, die kurzfristig 2.000+ Leads für 30.000\,\euro{} liefern würde. Aus Senior-Level-Sicht: Diese 2.000 Leads würden den Vertrieb mit 8 Personen überlasten, davon wären nur 20\,--\,40 ernsthaft \textendash{} und der CRO würde sich monatelang mit Vertriebskonflikten beschäftigen. Stattdessen: 30.000\,\euro{} in LinkedIn + Webinar, die 150 hochqualifizierte Leads liefern.

\begin{merkbox}[Senior-Level-Hinweis]
\emph{Performance-Marketing-Architektur ist eine Investitionsentscheidung, keine Werbeplanung.} Wer sie als Werbung versteht, optimiert auf Klicks. Wer sie als Investition versteht, optimiert auf Pipeline und Deckungsbeitrag \textendash{} und kann jede Kampagne der GF gegenüber wirtschaftlich verantworten.
\end{merkbox}
\end{loesungbox}

\clearpage

\section*{Abschluss}
\thispagestyle{plain}

\noindent Diskussionsbasis. Vergleichen Sie Ihre Annahmen, Ihre Zielkonflikte und Ihre Entscheidungen.

\medskip
\begin{infobox}[Nächster Schritt]
Coaching-Frage: Welche Kampagne hätten Sie kurzfristig gerne gestartet \textendash{} und welche Argumente sprechen dauerhaft gegen sie, auch wenn das Quartals-Lead-Ziel sie nahelegt?
\end{infobox}

\vspace{1cm}
\noindent{\color{lpAccent}\rule{\linewidth}{0.6pt}}\\[4pt]
{\footnotesize\sffamily\color{lpGray}Lernplatt \textperiodcentered{} by Can Siebert \textperiodcentered{} Kurs 22 (Bo 143) \textperiodcentered{} Tag 9 \textperiodcentered{} Lösungsheft \textperiodcentered{} \textcopyright{} 2026}

\end{document}
